\chapter{Probability}\label{chap:probability}

\section{Bounded random variables + concentration}\label{sec:concentration}

\begin{defn}[Bounded random variable]
\label{def:bounded-rv}
  \lean{ProbabilityTheory.BoundedRandomVariable}
  \leanok
$f : \Omega \to \R$ is \emph{bounded} in $[a, b]$ under $\mu$ iff
$f$ is measurable and $f(\omega) \in [a, b]$ for $\mu$-almost every
$\omega$.
\end{defn}

\begin{thm}[Chebyshev majority bound]
\label{thm:chebyshev-majority}
  \lean{ProbabilityTheory.chebyshev\_majority\_bound}
  \leanok
  \uses{def:bounded-rv}
Let $X_1, \dots, X_k$ be independent random variables with each
$\mu(A_j) \ge 2/3$ where $A_j = \{ X_j = \text{correct} \}$. Then for
$k \ge 9 / \delta$, the probability that a majority of the $X_j$ is
correct is at least $1 - \delta$. Proof via Popoviciu's variance
bound + Chebyshev's inequality \cite[\S2]{boucheron2013concentration}.
\end{thm}

\section{Exchangeability and double-sample}\label{sec:exchangeability}

\begin{defn}[Sample bundle]
\label{def:sample-bundle}
  \lean{ProbabilityTheory.ExchangeableSample}
  \leanok
A bundle $\langle m, \mu, \mathrm{IsProbabilityMeasure}\,\mu \rangle$
where $m \in \N$ is the sample size and $\mu$ is a probability measure
on the sample space. Used with the product measure
$\mu^{\otimes 2m}$ as the natural exchangeable double-sample domain.
\end{defn}

\begin{defn}[Valid split]
\label{def:valid-split}
  \lean{ProbabilityTheory.ValidSplit}
  \leanok
$s : \mathrm{Fin}(2m) \to \mathrm{Bool}$ is a \emph{valid split} iff
exactly $m$ of its values are \texttt{true}. There are
$\binom{2m}{m}$ valid splits (Vapnik-Chervonenkis double-sample
\cite{vapnik1971}).
\end{defn}

\begin{defn}[Uniform split measure]
\label{def:split-measure}
  \lean{ProbabilityTheory.SplitMeasure}
  \leanok
  \uses{def:valid-split}
The uniform probability measure on the set of valid splits. This is
the symmetrization measure for the double-sample trick.
\end{defn}

\section{Finite-type PMF}\label{sec:fintype-pmf}

\begin{defn}[Real-valued fintype PMF]
\label{def:fintype-pmf}
  \lean{ProbabilityTheory.FintypePMF}
  \leanok
For a fintype $\alpha$, a \texttt{FintypePMF}~$\alpha$ is a pair
$\langle p, h_{\ge 0}, h_{\sum = 1} \rangle$ where $p : \alpha \to \R$,
$p \ge 0$ pointwise, and $\sum_{a : \alpha} p(a) = 1$. This is an
$\R$-valued, \texttt{linarith}-compatible specialization of Mathlib's
$\R_{\ge 0}^\infty$-valued \texttt{PMF} \cite[\S4]{boucheron2013concentration}.
\end{defn}

\begin{thm}[Bridge to Mathlib PMF]
\label{thm:fintype-pmf-bridge}
  \lean{ProbabilityTheory.FintypePMF.toPMF\_apply}
  \leanok
  \uses{def:fintype-pmf}
The map $\mathrm{toPMF} : \mathrm{FintypePMF}\,\alpha \to \mathrm{PMF}\,\alpha$
satisfies $(\mathrm{toPMF}\,p)(a) = \mathrm{ENNReal.ofReal}(p.\mathrm{prob}\,a)$,
which is lossless on valid FintypePMF inputs.
\end{thm}

\begin{defn}[$L^1$ TV distance]
\label{def:tv-distance-fintype}
  \lean{ProbabilityTheory.FintypePMF.tvDistance}
  \leanok
$\mathrm{tvDistance}\,p\,q \;:=\; \sum_{a : \alpha} | p(a) - q(a) |$.
This is the \emph{full $L^1$} convention, not the half-sup TV of
Definition~\ref{def:tv-dist-real}; on discrete supports they differ
by a factor of $2$.
\end{defn}

\begin{thm}[Basic properties]
\label{thm:tv-distance-props}
  \lean{ProbabilityTheory.FintypePMF.tvDistance\_nonneg, ProbabilityTheory.FintypePMF.tvDistance\_comm, ProbabilityTheory.FintypePMF.tvDistance\_self}
  \leanok
  \uses{def:tv-distance-fintype}
$\mathrm{tvDistance}$ is non-negative, symmetric, and $0$ on
$p = q$.
\end{thm}

\begin{thm}[Transfer principle]
\label{thm:tv-transfer}
  \lean{ProbabilityTheory.FintypePMF.expectation\_approx\_of\_tv}
  \leanok
  \uses{def:tv-distance-fintype}
For Boolean-valued $f$,
$\bigl| \E_p[f] - \E_q[f] \bigr| \;\le\; \mathrm{tvDistance}(p, q)$.
In the half-sup TV convention this would read $\le 2 \cdot \TV(p, q)$.
\end{thm}
